\section{Limitations and threats to validity}
\label{sec:limitations}

An independent audit of the code and the frozen files, run after every result was frozen,
produced the findings below; the full list, with evidence and the disposition of each, is in
\file{docs/AUDIT.md}. No finding changed a frozen file. Where a finding could be addressed from
committed data without a new rollout, it is addressed in this report and labelled.

\paragraph{Few seeds, and verdicts that depend on the method.} Each family has
\rf{proto.nseeds} training seeds. The seed-level sensitivity analysis
(Tables~\ref{tab:sens-main} and~\ref{tab:sens-rob}) changes
\rf{sens.l1.changed} of \rf{sens.l1.n} Level~1 verdicts, \rf{sens.rob.changed} of
\rf{sens.rob.n} robustness and ablation verdicts, and \rf{sens.ext.changed} of
\rf{sens.ext.n} verdicts of the two extensions; it is undefined for
\rf{sens.l1.undefined}, \rf{sens.rob.undefined} and \rf{sens.ext.undefined} of them because
the seeds do not vary. The separated Can H2 and H3 results, and several robustness
separations, are among the verdicts that change. The audit also found two verdicts that change
with the bootstrap's random seed. Of the \rf{sens.all.n} differences in the sensitivity
tables (the declared ones and one context row), \rf{sens.all.boundary} have a frozen interval
with a bound at exactly zero or of zero width; with \rf{proto.nstates} binary outcomes per seed the bootstrap distribution is
discrete, so these verdicts rest on ties. A policy that succeeds on all \rf{proto.nstates}
starts has a Wilson interval of [\rf{wilson.full.lo}, \rf{wilson.full.hi}]\,\%, so near the
ceiling the test separates only large differences.

\paragraph{Multiplicity.} No interval is adjusted for multiplicity. Each Level~1 task has
\rf{proto.ncomparisons} success comparisons and \rf{l1.lift.nlen} length comparisons, and
the robustness analyses report \rf{rob.lift.contrasts} (Lift) and \rf{rob.can.contrasts}
(Can) differences and contrasts. A Bonferroni correction within each analysis would remove several separations, among them
Can H2; the audit (\file{docs/AUDIT.md}, B-1 and B-6) gives the corrected intervals and the
count.

\paragraph{The diffusion sampler bounds its actions.} The sampler clips the predicted clean
action chunk to $[-1, 1]$ in min--max normalized space, so every diffusion action lies within
the range of the training actions on each dimension. The BC heads are clipped only to the
controller limits, which are wider. The audit found that in BC episodes that fail under perturbation most motion actions lie
outside the training range, while diffusion's never can. This is a correlation, not a shown
mechanism: the out-of-range actions begin only after a grasp is missed, BC-MLP leaves the
range on many base starts that it still solves, and BC-RNN rarely leaves it yet loses half
its success under the Lift arm offset. This bound is part of what H2-c compares: diffusion and BC-MLP-chunk differ in the
generative head \emph{and} in this output bound, so neither the Level~1 nor the robustness
differences between them isolate the generative model. Clamping BC outputs to the training
range would test this without retraining; it was not run in this report, because it would be
a new evaluation after the freeze.

\paragraph{Selection and its bias.} Validation success, which chose among \rf{train.K}
checkpoints and \rf{train.trials} trials, is optimistic: Table~\ref{tab:selection} shows it
beside success on the fresh base starts and on the held-out starts (for example BC-MLP on
Can: \rf{sel.can.bcmlp.val}\,\% on validation, \rf{l1.can.bcmlp.succ}\,\% held out).
Of the \rf{l1.bc.nruns} BC runs, \rf{l1.bc.first-ckpt} selected their first checkpoint, and two
diffusion runs their last, so the checkpoint grid may have been too coarse for BC and the
budget too short for those diffusion runs. In the RGB and Level~2 studies most checkpoints succeeded on
every validation start, so selection there was decided by tiny differences in episode length
and carries almost no information about the quantities those studies test; one RGB diffusion
run and one point-cloud run selected their final checkpoint. Each family had only
\rf{train.trials} tuning trials, and the extensions none, so no family is tuned to its best. Before the Lift
search space was committed, a short pipeline rehearsal (a few hundred training steps) had
evaluated the default and the higher-learning-rate configurations of BC-MLP and diffusion on
the validation suite; the higher learning rate was later chosen for Lift diffusion. The trial
list was the same for every family, but BC-RNN and BC-MLP-chunk had no such early look.

\paragraph{One rollout per start.} With $R = \rf{proto.repeats}$, a start's difficulty and the
diffusion noise drawn on it are confounded, so diffusion's interval mixes policy stochasticity
with start difficulty, unlike the deterministic BC families (KI-028).

\paragraph{Reuse of starts across studies.} The same \rf{proto.nstates} Lift held-out starts
served the Level~1 test, the RGB test and the Level~2 held-out check, and the robustness base
and cube-shift starts served both H4 and H5. No model, checkpoint or configuration was ever
selected on them. However, the three studies were designed in sequence and later designs were
informed by earlier outcomes on these starts: H5's primary set was chosen after the Level~1
robustness and RGB held-out results were known. The Level~2 held-out comparison, with both
arms at the ceiling, could not have favoured the point-cloud arm and is a replication-type
check, not a test.

\paragraph{Scale of the diffusion policy.} The diffusion family was matched to BC in
parameters (about a million) and in its optimizer recipe. robomimic's default U-Net is wider
and deeper and is trained with a warm-up cosine schedule, so these results speak to
diffusion at this scale and recipe only.

\paragraph{In-distribution starts.} The evaluation starts come from the same reset
distribution as the demonstrations; on Lift each start lies within a few millimetres of a
demonstration start (KI-011), so Lift held-out success measures in-distribution performance.
Open-loop replay of the Can demonstrations diverges at contacts by several centimetres
(KI-004), so the simulator reproduces the recorded dynamics only approximately.

\paragraph{Privilege, features and platform.} Every Level~1 result uses exact object poses.
The Can input recomputes the gripper-to-object vector from the current step, unlike
robomimic's stored feature, so Can numbers are not comparable with robomimic's. The data are
the robosuite~v1.5 repackaging of the robomimic datasets, so no number here should be compared
with the original benchmark. All rollouts ran in simulation on one platform. Our success
predicate is evaluated on the current step's kinematics, while robosuite's own predicate reads
the previous substep; in a small fraction of episodes ours fires one control step earlier,
which shifts first-success steps but not success. The point cloud uses noise-free depth, exact
camera calibration and a crop box whose lower edge removes the table at its known height.

\paragraph{Declared quantities not reported, and quantities not declared.} The protocol's
secondary metrics (gripper sign changes, action smoothness) were computed per episode but
never aggregated by the frozen analysis; Table~\ref{tab:secondary} reports them post hoc. The diffusion-specific plan-diversity metric was measured only in a pilot and is not
reported, and the point-cloud throughput listed in the Level~2 freeze is not in its results.
Validation success at every checkpoint, declared for the RGB study, and the point-cloud
per-checkpoint rates listed in the Level~2 freeze are in the committed selection files but
not in this report. The H4 analysis
compares four of the six family pairs, a choice fixed in code before any robustness rollout.
The optional DiT and GMM-BC families and the Square task were not run, to stay within the
project's time and compute budget; the declared hypotheses did not depend on them.

\paragraph{Analysis code, records and cost figures.} The Level~1 analysis script was
hash-frozen; the robustness, RGB and Level~2 analysis scripts were not, and each received a
small fix after its data existed. The fixes touched input reading and training-time accounting
only, not pairing, resampling or which evaluations enter. The validation summaries of one RGB
training run were recorded from a working tree with uncommitted changes; the next commit
changed only documentation, logs and reports, and the frozen code hashes match. The Level~2 analysis misreported one training time (KI-049), and one RGB run's
hours leave out a preempted attempt (Table~\ref{tab:cost}). Billed amounts are Modal's usage
report at list price before credits, not invoices; one early approval allowed one GPU probe
and two probe launches ran before the owner amended it.

\paragraph{Scope of the claims.} The comparison covers two tasks, privileged state for
Level~1, \rf{data.lift.train} training demonstrations per task from one operator channel, one training budget and one
platform. Within that scope the evidence supports an advantage of BC-RNN over BC-MLP on Lift under
both interval methods and on Can under the preregistered bootstrap only, gives no evidence
that a diffusion head adds to a deterministic head with the same horizons, and shows that our point-cloud policy generalizes worse than our RGB
policy to unseen cube positions. It supports no claim about real robots, other tasks, larger
datasets or DP3 itself.
